\documentclass[12pt,a4paper]{amsart}

\usepackage{amsmath,amssymb,amsthm}
\usepackage{mathtools}
\usepackage{enumitem}
\usepackage{hyperref}

\newtheorem{theorem}{Theorem}[section]
\newtheorem{lemma}[theorem]{Lemma}
\newtheorem{proposition}[theorem]{Proposition}
\newtheorem{corollary}[theorem]{Corollary}
\theoremstyle{definition}
\newtheorem{definition}[theorem]{Definition}
\newtheorem{remark}[theorem]{Remark}
\newtheorem{problem}[theorem]{Open Problem}

\newcommand{\Rig}{\textup{[Rig]}}
\newcommand{\Num}{\textup{[Num]}}
\newcommand{\Warn}{\textup{[Warning]}}
\newcommand{\Open}{\textup{[Open]}}

\begin{document}

\title{Fold Amplitude with Singular Weight:\\
  Asymptotic Phase Diagram and CFU Reduction}

\author{[Author]}

\begin{abstract}
We study the oscillatory integral
\[
  \Phi_\alpha(u)
  :=\int_0^\infty
  \beta^{-1/2}
  \chi\!\left(\frac{\beta}{|u|^\alpha}\right)
  e^{i(-|u|\beta + c_3\beta^3)}\,d\beta,
  \quad u<0,\; c_3>0.
\]
The interaction between the cubic phase and the boundary singularity
$\beta^{-1/2}$ induces a scaling transition at $\alpha=\tfrac{1}{2}$,
separating a classical Erd\'elyi regime from a CFU-type
coalescence regime.
We prove the uniform bound
$|\Phi_\alpha(u)| = O(|u|^{-1/2})$ for all $\alpha\ge 0$.
For $\alpha<\tfrac{1}{2}$ this decay is sharp
($\Theta(|u|^{-1/2})$) with leading constant $\sqrt{\pi}$.
For $\alpha\ge\tfrac{1}{2}$, no limit of the rescaled amplitude
$|u|^{1/2}|\Phi_\alpha(u)|$ exists:
$\limsup=2\sqrt{\pi}$ and $\liminf=0$
(along the continuous variable $u\to-\infty$),
due to persistent interference between
endpoint and interior saddle contributions.
Earlier claimed exponents $|u|^{-1/6}$ and $|u|^{-1/3}$
are shown to be artifacts of an incorrect scaling regime.
\end{abstract}

\maketitle

\begin{theorem}[Universal decay]
\label{thm:decay}
Let $\alpha\ge 0$, $c_3>0$, $u\to-\infty$.
\begin{enumerate}[\rm(i)]
\item \textup{(Uniform bound)}
  \[
    |\Phi_\alpha(u)| = O(|u|^{-1/2}).
  \]
\item \textup{(Sharp regime, $\alpha<\tfrac{1}{2}$)}
  \[
    |\Phi_\alpha(u)| = \Theta(|u|^{-1/2}),
    \quad\text{with leading constant }\sqrt{\pi}.
  \]
\end{enumerate}
\end{theorem}

\begin{theorem}[Non-existence of limit amplitude]
\label{thm:nolimit}
For $\alpha\ge\tfrac{1}{2}$, $c_3>0$,
the following holds for the continuous variable $u\to-\infty$:
\[
  \limsup_{u\to-\infty}|u|^{1/2}|\Phi_\alpha(u)| = 2\sqrt{\pi},
  \qquad
  \liminf_{u\to-\infty}|u|^{1/2}|\Phi_\alpha(u)| = 0.
\]
In particular, no constant $C$ satisfies
$|\Phi_\alpha(u)|\sim C|u|^{-1/2}$.
\end{theorem}

\section{Phase Boundary}
\label{sec:boundary}

\begin{definition}[Saddle and boundaries]
\label{def:boundary}
The saddle of $\phi(\beta)=-|u|\beta+c_3\beta^3$
sits at $\beta_0(u)=\sqrt{|u|/(3c_3)}$;
the support has effective radius $R(u)=|u|^\alpha$.

\medskip
\noindent\textbf{Asymptotic scaling boundary:}
$\alpha=\tfrac{1}{2}$.
\begin{itemize}
\item $\alpha<\tfrac{1}{2}$:
  $\beta_0/R\to\infty$;
  saddle exits support;
  Erd\'elyi applies.
\item $\alpha>\tfrac{1}{2}$:
  $\beta_0/R\to 0$;
  saddle coalesces with
  singularity at $\beta=0$;
  CFU-type rescaling is natural.
\item $\alpha=\tfrac{1}{2}$:
  $\beta_0/R=(3c_3)^{-1/2}$ fixed;
  structural CFU boundary.
\end{itemize}

\medskip
\noindent\textbf{Finite-$|u|$ crossover curve:}
$\beta_0(u)=R(u)$ at
$|u|_{\rm cross}=(3c_3)^{1/(1-2\alpha)}$
(for $\alpha<\tfrac{1}{2}$).
This is a finite-size effect,
not the asymptotic boundary.
\end{definition}

\section{Region I: Erd\'elyi}
\label{sec:erdelyi}

\begin{theorem}[Sharp Erd\'elyi \Rig]
\label{thm:erdelyi}
For $\alpha=0$ and $c_3>0$:
\[
  \Phi_0(u)
  =\sqrt{\pi}\,e^{-i\pi/4}
  |u|^{-1/2}+O(c_3|u|^{-3/2}).
\]
\end{theorem}

\begin{corollary}[Ces\`aro sharp \Rig]
$\limsup_{u\to-\infty}
|u|^{1/2}|\Phi_0(u)|
=\sqrt{\pi}>0$;
hence $|F_0(U)|=\Theta(U^{1/2})$.
\end{corollary}

\section{Region III: CFU Reduction}
\label{sec:cfu}

\begin{theorem}[CFU Reduction \Rig]
\label{thm:cfu}
For $\alpha\ge\tfrac{1}{2}$,
the substitution $\beta=\beta_0(u)\cdot x$
with $\beta_0=\sqrt{|u|/(3c_3)}$
admits the asymptotic representation
\begin{equation}\label{eq:cfu}
  \Phi_\alpha(u)
  = \beta_0^{1/2}\cdot I(\lambda)
  + O(|u|^{-1}),
  \qquad
  \lambda
  := \frac{|u|^{3/2}}{(3c_3)^{1/2}},
\end{equation}
where
$I(\lambda)
:= \int_0^{X(\lambda)}
t^{-1/2}\,e^{i\lambda(t^3/3-t)}\,dt$
and $X(\lambda)=|u|^{\alpha-1/2}$.
For $\alpha>\tfrac{1}{2}$, $X(\lambda)\to\infty$
and the integral can be replaced asymptotically
by its extension to $[0,\infty)$
with a remainder $O(\lambda^{-N})$ for any $N$
(formally, by repeated integration by parts
on $[X(\lambda),\infty)$).
For $\alpha=\tfrac{1}{2}$, $X(\lambda)=1$ is fixed;
the saddle $t=1$ lies at the boundary of integration,
yielding a half-saddle contribution
in the sense of endpoint stationary phase asymptotics
(factor $\tfrac{1}{2}$ relative to an interior saddle~\cite{Olver1974});
the order $O(\lambda^{-1/2})$ is unchanged.

\medskip
\noindent\textbf{Jacobian:}
\[
  \beta^{-1/2}\,d\beta
  = \beta_0^{1/2}\cdot x^{-1/2}\,dx,
  \qquad
  \beta_0^{1/2}
  = \Bigl(\frac{|u|}{3c_3}\Bigr)^{1/4}
  = |u|^{1/4}(3c_3)^{-1/4}.
\]

\medskip
\noindent\textbf{Phase:}
\[
  -|u|\beta+c_3\beta^3
  = |u|\beta_0\bigl(-x+\tfrac{1}{3}x^3\bigr)
  = \lambda\bigl(\tfrac{x^3}{3}-x\bigr),
\]
confirming $\lambda=|u|^{3/2}(3c_3)^{-1/2}$.

\medskip
\noindent\textbf{Role of CFU rescaling.}
For $\alpha\ge\tfrac{1}{2}$, classical Erd\'elyi-type expansions
are not uniform near the coalescence regime:
the phase and amplitude singularity scale on the same asymptotic
order, so a single Erd\'elyi-type expansion does not
capture both contributions uniformly as $u\to-\infty$.
The CFU-type substitution $\beta\mapsto\beta_0 x$ is natural here:
it simultaneously resolves the large parameter
to $\lambda\sim|u|^{3/2}$ and separates endpoint
from saddle contributions.

\medskip
\noindent\textbf{Two-term asymptotics of $I(\lambda)$.}
For $\alpha>\tfrac{1}{2}$ and large $\lambda$,
$I(\lambda)$ admits the asymptotic decomposition
into endpoint and saddle contributions:
\begin{itemize}
\item \emph{Endpoint} $t=0$
  (Erd\'elyi, weight $t^{-1/2}$):
  \[
    I_{\rm end}(\lambda)
    = \sqrt{\pi}\,e^{-i\pi/4}
    \lambda^{-1/2}
    + O(\lambda^{-3/2}).
  \]
\item \emph{Saddle} $t=1$
  (stationary phase,
  phase $\lambda(-\tfrac{2}{3})$,
  amplitude $1^{-1/2}=1$,
  curvature $\phi''(1)=2$):
  \[
    I_{\rm sad}(\lambda)
    = \sqrt{\tfrac{\pi}{\lambda}}
    \,e^{-2i\lambda/3+i\pi/4}
    + O(\lambda^{-3/2}).
  \]
\end{itemize}
Both are $O(\lambda^{-1/2})$.
Summing the endpoint and saddle contributions:
\begin{equation}\label{eq:twoterm}
  I(\lambda)
  = \sqrt{\pi}\,\lambda^{-1/2}
  \Bigl[
    e^{-i\pi/4}
    + e^{-2i\lambda/3+i\pi/4}
  \Bigr]
  + O(\lambda^{-3/2}),
\end{equation}
\begin{equation}\label{eq:Ibound}
  |I(\lambda)|
  \le 2\sqrt{\pi}\,\lambda^{-1/2}.
\end{equation}

\begin{remark}[Boundary saddle at $\alpha=\tfrac{1}{2}$]
\label{rem:halfsaddle}
For $\alpha=\tfrac{1}{2}$, the upper limit $X(\lambda)=1$
coincides with the saddle $t=1$.
The stationary-phase contribution from a saddle
at the endpoint of integration is halved relative
to an interior saddle contribution,
in the sense of endpoint stationary phase
asymptotics~\cite{Olver1974}.
This modifies only the constant in~\eqref{eq:twoterm},
not the decay rate. In particular,
$|I(\lambda)|=O(\lambda^{-1/2})$ still holds,
and the final bound~\eqref{eq:final} is unaffected.
\end{remark}

\medskip
\noindent\textbf{Final bound.}
Combining~\eqref{eq:cfu}
with $\beta_0^{1/2}=|u|^{1/4}(3c_3)^{-1/4}$
and $|I(\lambda)|=O(\lambda^{-1/2})
=O(|u|^{-3/4})$:
\begin{equation}\label{eq:final}
  \boxed{
  |\Phi_\alpha(u)|
  = O\bigl(|u|^{1/4}\cdot|u|^{-3/4}\bigr)
  = O\bigl(|u|^{-1/2}\bigr).
  }
\end{equation}

\begin{center}
\fbox{\parbox{0.85\linewidth}{
  \textbf{Structural result:}
  In the presence of the boundary
  singularity $\beta^{-1/2}$,
  the saddle contributes at the
  same order $O(\lambda^{-1/2})$
  as the endpoint.
  It does not improve the decay rate;
  it only modulates the amplitude
  through the oscillating factor
  $e^{-2i\lambda/3}$ in~\eqref{eq:twoterm}.
}}
\end{center}
\end{theorem}

\begin{remark}[No limit of the rescaled amplitude,
  Region~III \Warn]
\label{rem:nolimit}

\noindent\textbf{Coefficient.}
From~\eqref{eq:cfu} we have
\[
  \beta_0^{1/2}\lambda^{-1/2}
  =
  \Bigl(\frac{|u|}{3c_3}\Bigr)^{1/4}
  \cdot
  \Bigl(\frac{(3c_3)^{1/2}}{|u|^{3/2}}\Bigr)^{1/2}
  =
  |u|^{-1/2},
\]
so all $c_3$-dependence cancels in the leading term;
this reflects a scaling invariance of the cubic phase.

\medskip
Combining with~\eqref{eq:twoterm}:
\begin{equation}\label{eq:rescaled}
  |u|^{1/2}|\Phi_\alpha(u)|
  = \sqrt{\pi}
  \cdot
  \Bigl|e^{-i\pi/4}
  +e^{-2i\lambda/3+i\pi/4}\Bigr|
  + O(|u|^{-1}),
\end{equation}
where $\lambda=|u|^{3/2}(3c_3)^{-1/2}\to\infty$.
The factor $|e^{-i\pi/4}+e^{-2i\lambda/3+i\pi/4}|$
oscillates between $0$ and $2$
without converging;
this oscillation is the sole source
of non-convergence of the rescaled amplitude.
This is distinct from
``no power law'': the exponent
$-1/2$ is correct;
it is the \emph{coefficient}
that fails to converge.
\end{remark}

\begin{remark}[Proof of Theorem~\ref{thm:nolimit}]
\label{rem:extremal}
From~\eqref{eq:rescaled}, writing
$\theta:=\tfrac{2}{3}\lambda(u)=\tfrac{2}{3(3c_3)^{1/2}}|u|^{3/2}$:
\[
  |u|^{1/2}|\Phi_\alpha(u)|
  = \sqrt{\pi}\,
  \bigl|1 + e^{-i(\theta-\pi/2)}\bigr|
  + O(|u|^{-1}).
\]

\noindent\textbf{Extremal values of the main term.}
The factor $|1+e^{-i(\theta-\pi/2)}|$
equals $2$ when $\theta\equiv\tfrac{\pi}{2}\pmod{2\pi}$
(constructive interference)
and $0$ when $\theta\equiv\tfrac{3\pi}{2}\pmod{2\pi}$
(destructive interference).
Since $f(u) = A(u) + o(1)$ implies
$\limsup f = \limsup A$ and $\liminf f = \liminf A$,
the $O(|u|^{-1}) = o(1)$ remainder does not affect
the $\limsup$ and $\liminf$ of~\eqref{eq:rescaled}.

\medskip
\noindent\textbf{Surjectivity of the phase.}
The function $\theta:(-\infty,u_0]\to(0,\infty)$
is continuous and strictly increasing,
hence its image is an unbounded interval;
in particular $\theta$ traverses every interval
$[k\cdot 2\pi,(k+1)\cdot 2\pi]$ for all large $k$,
so every phase value in $[0,2\pi)$ is
\emph{surjectively attained} by $\theta(u)\bmod 2\pi$.
No density or number-theoretic argument is required.

\medskip
\noindent\textbf{Explicit sequences.}
For any target $\varphi\in[0,2\pi)$,
the exact inversion $\theta(u_n)=2\pi n+\varphi$ gives
\[
  |u_n^{(\varphi)}|
  = \Bigl(\tfrac{3}{2}(3c_3)^{1/2}(2\pi n+\varphi)\Bigr)^{2/3},
  \qquad
  u_n^{(\varphi)}\to-\infty.
\]
For $\varphi=\tfrac{\pi}{2}$: the leading term in~\eqref{eq:rescaled}
equals $2\sqrt{\pi}$, up to an $o(1)$ correction,
so $|u_n|^{1/2}|\Phi_\alpha(u_n)|\to 2\sqrt{\pi}$.
For $\varphi=\tfrac{3\pi}{2}$: the leading term equals $0$,
up to an $o(1)$ correction,
so $|u_n|^{1/2}|\Phi_\alpha(u_n)|\to 0$.

\medskip
\noindent\textbf{Conclusion.}
$\limsup$ and $\liminf$ are taken with respect
to the continuous variable $u\to-\infty$.
The sequences above yield
Theorem~\ref{thm:nolimit}. \qed
\end{remark}

\begin{remark}[Retraction of
  $-1/6$ and $-1/3$ \Warn]
\label{rem:retract}
The claims $|\Phi_\alpha|\lesssim|u|^{-1/6}$
and $\sim|u|^{-1/3}$ arose from:
\begin{enumerate}[{\rm(i)}]
\item Using the standard Airy
  parameter $\lambda\sim|u|^{2/3}$
  (valid without singular weight).
  Here $\lambda=|u|^{3/2}(3c_3)^{-1/2}$.
\item Ignoring the Jacobian
  $\beta_0^{1/2}=|u|^{1/4}(3c_3)^{-1/4}$
  which upgrades $\lambda^{-1/2}
  \sim|u|^{-3/4}$ to
  the correct $|u|^{-1/2}$.
\end{enumerate}
\end{remark}

\section{Numerical Evidence}
\label{sec:num}

\begin{proposition}[\Num]
Fitted exponents $\gamma$ from
non-uniform log-log regression of $|\Phi_\alpha(u)|$
(inherently unstable due to phase modulation),
$c_3=R_0=1$:
\begin{center}
\begin{tabular}{c|c|c|c|c}
$\alpha$ & region
  & $\gamma([16,50])$
  & $\gamma([50,200])$
  & $\gamma([200,1000])$
\\\hline
$0.00$ & I & $-0.502$ & $-0.500$ & $-0.500$ \\
$0.25$ & I/II & $-0.391$ & $-0.489$ & $+0.551$ \\
$0.50$ & III & $-0.299$ & $-1.235$ & $-0.071$ \\
$0.75$ & III & $-0.337$ & $-0.053$ & $+0.334$ \\
$1.00$ & III & $-0.558$ & $+0.175$ & $-0.343$
\end{tabular}
\end{center}
Only $\alpha=0$ shows convergence
of the fitted exponent.
For $\alpha\ge\tfrac{1}{2}$,
the wildly varying fits reflect non-uniform log-log
regression in the presence of phase modulation:
the oscillating factor $|e^{-i\pi/4}+e^{-2i\lambda/3+i\pi/4}|$
can be near $0$ or near $2$ on any given interval,
giving apparent exponents far from $-1/2$.
The extreme value $\gamma\approx-1.235$
for $\alpha=0.5$ on $[50,200]$ reflects
proximity to a near-zero of the amplitude factor;
these regression values are heuristic only.
\end{proposition}

\section{The PSWF Open Problem}
\label{sec:PSWF}

\begin{problem}[\Open]
\label{prob:PSWF}
For the fold integral in Paper~XV,
determine $\alpha$:
\begin{enumerate}[{\rm(i)}]
\item If $\alpha<\tfrac{1}{2}$:
  $|\Phi|=\Theta(|u|^{-1/2})$;
  Ces\`aro bound sharp.
\item If $\alpha\ge\tfrac{1}{2}$:
  $|\Phi|=O(|u|^{-1/2})$;
  no limit of the rescaled amplitude;
  Ces\`aro sharpness depends
  on cancellation in $F(U)$.
\end{enumerate}
$\alpha$ is determined by the PSWF spectral expansion.
This corresponds to the $\mu=-\tfrac{1}{2}$ slice
of a more general $(\mu,\alpha)$ phase diagram,
where the weight $\beta^{-1/2}$ is replaced by $\beta^\mu$;
the full classification of regimes in $(\mu,\alpha)$
is deferred to a subsequent work.
\end{problem}

\begin{thebibliography}{9}
\bibitem{Erdelyi1956}
A.~Erd\'elyi,
\textit{Asymptotic Expansions},
Dover, 1956.
\bibitem{Olver1974}
F.~W.~J.~Olver,
\textit{Asymptotics and
Special Functions},
Academic Press, 1974.
\bibitem{CFU1957}
C.~Chester, B.~Friedman,
F.~Ursell,
\textit{An extension of the
method of steepest descents},
Proc.\ Cambridge Philos.\ Soc.\
\textbf{53} (1957), 599--611.
\end{thebibliography}

\end{document}
